Para SKYFILMS, la gestión de riesgos de seguridad operacional es un pilar fundamental para garantizar la seguridad de nuestras operaciones con UAS. Reconocemos que el entorno aeronáutico presenta peligros inherentes, y nuestra capacidad para identificarlos y gestionar los riesgos asociados es crucial para prevenir incidentes y accidentes. Este proceso se fundamenta en un enfoque sistemático y formal que involucra la identificación de peligros, la evaluación, mitigación y aceptación de los riesgos de seguridad operacional.\\

La efectividad de esta gestión reside en la participación activa y el compromiso de cada individuo dentro de nuestra organización, desde la alta dirección hasta cada piloto. Para ello, SKYFILMS ha establecido un proceso formal y continuo para la identificación de peligros, combinando métodos reactivos, proactivos y predictivos, asegurando una cobertura integral de los posibles riesgos.\\

\section{Identificación de Peligros}
En SKYFILMS, la identificación de peligros es el primer y crucial paso en nuestra gestión de la seguridad operacional, entendiendo un peligro como una condición u objeto que podría provocar un incidente o accidente o contribuir al mismo. Nos centramos en una amplia gama de factores que pueden impactar nuestras operaciones con UAS, garantizando así el cumplimiento de los principios del SMS. Para ello, se consideran los siguientes criterios principales:

\begin{enumerate}
    \item \textbf{Peligros al personal operativo:} Condiciones en la operación que afectan directamente la integridad física y mental de nuestro equipo.
    \item \textbf{Peligros a personas ajenas a la operación:} Potenciales amenazas para la seguridad de terceros en tierra o en el aire.
    \item \textbf{Peligros en tierra (obstáculos naturales y artificiales):} Identificación de estructuras, terrenos, flora o fauna que puedan interferir con la operación.
    \item \textbf{Peligros en el aire (por fallas en la operación):} Riesgos derivados de posibles mal funcionamiento del equipo, errores humanos o condiciones operacionales adversas durante el vuelo.
    \item \textbf{Peligros estratégicos (por la interacción en espacios aéreos de aeronaves tripuladas):} Peligros asociados a la coexistencia con aeronaves tripuladas en el espacio aéreo, especialmente en zonas de operación compartidas.
    \item \textbf{Peligros por cambios en la normativa vigente aplicable:} La necesidad de adaptar constantemente nuestros procedimientos a las modificaciones en el RAC 100 y RAC 219, y otras regulaciones relevantes.
    \item \textbf{Peligros por la gestión del cambio:} Evaluación de los riesgos introducidos por cualquier modificación en nuestros procesos, equipos o personal, antes de su implementación.
\end{enumerate}

Para mantener nuestro Sistema de Gestión de Seguridad Operacional actualizado y en constante búsqueda de posibles peligros que aumenten el riesgo de las operaciones, SKYFILMS utiliza los métodos de identificación expuestos a continuación.

\subsection{Método Reactivo}
Este método SKYFILMS lo aplica para identificar peligros y factores contribuyentes a partir de eventos que ya han ocurrido. Esto nos permite aprender de la experiencia y evitar la recurrencia de sucesos no deseados. Las principales herramientas reactivas de SKYFILMS son:\\

\textbf{a) Sistema de Notificaciones de Seguridad Operacional:} Es la principal fuente de información reactiva. Fomentamos activamente una cultura de notificación, garantizando un tratamiento confidencial, protección de datos y no punitivo para los reportes voluntarios, esta información es almacenada en archivo digital con el objetivo de tener trazabilidad y tener presente la condición o peligro identificado para su retroalimentación y correcta gestión. Este sistema se divide en:

\begin{itemize}
    \item \textbf{Notificaciones Obligatorias (MOR):} Conforme a la regulación, SKYFILMS notificará de forma obligatoria a la Aerocivil (Grupo de Investigación de Accidentes, DIACC), dentro de las 48 horas siguientes (RAC 100.240), cualquier accidente o incidente grave. Así mismo, se reportarán los eventos de seguridad operacional definidos en la normativa vigente \textbf{(Anexo 4 - Listado Reportes Obligatorios M.O.R)}. Toda notificación obligatoria amerita una investigación formal por parte de la Junta de Revisión de Seguridad Operacional SRB.
    \item \textbf{Notificaciones Voluntarias:} Se incentiva a todo el personal y a terceros a reportar cualquier condición, acto o situación que perciban como un peligro potencial, incluso si no resultó en un incidente. Esto es crucial para identificar condiciones latentes y tendencias antes de que escalen a eventos mayores, siguiendo la "Teoría del Iceberg". Para facilitar este proceso, se utilizará un formato estandarizado.
\end{itemize}

\textbf{NOTA 1:} En cualquier caso de notificación, ya sea obligatoria o voluntaria se realizara retroalimentación presencial o virtual a la persona que realizo la notificación, la \textbf{Junta de Revisión de Seguridad Operacional SRB} es la encargada de analizar dicha notificación y la retroalimentación estará a cargo del \textbf{Gerente SMS}.\\

\textbf{NOTA 2:} En el \textbf{“Anexo 5 - Guía de Uso Formato de Notificación Peligros de Seguridad Operacional”} se explica el método de uso del formato de notificaciones.\\

\textbf{b) Investigación de Sucesos de Seguridad Operacional:} No todas las notificaciones requerirán una investigación a fondo. El \textbf{Gerente SMS}, en conjunto con la \textbf{Junta de Revisión de Seguridad Operacional (SRB)}, clasificará los reportes según el nivel de riesgo potencial para determinar cuáles ameritan una investigación formal. El objetivo de toda investigación no es encontrar culpables, sino determinar las causas raíz (factores contribuyentes y organizacionales) para implementar defensas efectivas. Para ello, SKYFILMS podrá utilizar técnicas de análisis reconocidas como:

\begin{itemize}
    \item Los 5 Porqués.
    \item Diagrama de Ishikawa (Causa-Efecto o Espina de Pescado).
    \item Análisis Bow-Tie (Corbatín).
\end{itemize}

Los hallazgos de cada investigación se documentarán y se utilizarán para actualizar el registro de peligros y los procesos de gestión de riesgos.

\subsection{Métodos Proactivo y Predictivo}
Estos métodos se enfocan en la anticipación, buscando peligros mediante el análisis de las operaciones actuales (proactivo) y la proyección de problemas futuros a partir de datos históricos (predictivo). SKYFILMS los implementa para crear defensas robustas antes de que los peligros se materialicen de la siguiente manera:\\

\textbf{a) Inspecciones y Auditorías de Seguridad Operacional:}

\begin{itemize}
    \item \textbf{Auditorías Internas:} Se realizarán auditorías planificadas a los procesos del SMS y del Manual de Operaciones (MO) para verificar el cumplimiento y la efectividad de los procedimientos.
    \item \textbf{Inspecciones Pre-vuelo:} Antes de cada vuelo, el piloto al mando realizará una inspección exhaustiva del UAS, la estación de control (GCS) y los equipos asociados.
    \item \textbf{Análisis de Seguridad en Sitio:} Previo al inicio de cada proyecto, el Jefe de Pilotos o un piloto designado realizará un análisis en el lugar de la operación para identificar peligros específicos del entorno: obstáculos (naturales y artificiales), infraestructura crítica, áreas de posible interferencia electromagnética, condiciones de acceso y la presencia de personas ajenas a la operación.
\end{itemize}

\textbf{b) Herramienta de Evaluación de Riesgo de Vuelo (FRAT):} Antes de cada operación, el piloto remoto deberá completar el formato \textbf{F-011 Evaluación del Riesgo Operacional (Anexo 6 - Guía de Uso FRAT)}. Esta herramienta calcula el riesgo presente de la operación, asignando una puntuación de 1 a 5 a 12 factores de riesgo (ej. relacionados con el piloto, el entorno, la complejidad de la misión). Un puntaje elevado alertará al equipo sobre la necesidad de implementar medidas de mitigación adicionales antes de iniciar el vuelo.

\subsection{Fuentes de Identificación de Peligros}
\begin{longtable}{|>{\RaggedRight\arraybackslash}p{3cm}|>{\RaggedRight\arraybackslash}p{6.2cm}|>{\RaggedRight\arraybackslash}p{4.4cm}|}
\hline
\rowcolor[gray]{0.9}\textbf{Enfoque} & \textbf{Fuentes} & \textbf{Responsable} \\ \hline
\endhead
Reactivo & Reportes obligatorios y voluntarios, investigación de eventos, discrepancias de mantenimiento, quejas de clientes & Gerente SMS \\ \hline
Proactivo & FRAT, inspecciones prevuelo, análisis del sitio, auditorías, lista de verificación de la misión, observación de operaciones, cambios en procesos & Jefe de Pilotos y Gerente SMS \\ \hline
Predictivo & Tendencias de SPI, registros de vuelo y mantenimiento (por ejemplo, desgaste de motores o baterías), análisis de fatiga y estrés, cambios en normas & Gerente SMS \\ \hline
Internas & Personal, aeronaves, instalaciones, procedimientos y software & SAG \\ \hline
Externas & Aerocivil, clientes, clima, tráfico aéreo, proveedores, otros operadores & Gerente SMS \\ \hline
\end{longtable}

\textbf{Factores humanos y organizacionales.} Se miran la fatiga, el estrés, la presión del cliente, la carga de trabajo y la comunicación como fuente de peligros y se registran en el Anexo 7.

\subsection{Inicio de una Investigación y Perfil del Investigador}
\begin{itemize}
    \item \textbf{Cuándo se investiga:} siempre en eventos del Anexo 4 y en cualquier evento con riesgo rojo o amarillo alto, daño a la aeronave o repetición de un evento. En los demás casos, el Gerente SMS decide con un análisis inicial.
    \item \textbf{Quién investiga:} el Gerente SMS, que nombra por escrito al investigador designado (puede ser él mismo). Si el evento le involucra, el Gerente General nombra otra persona.
    \item \textbf{Perfil:} formación en SMS o investigación de eventos (curso del Anexo 12), conocimiento de la operación UAS, independencia frente a las personas involucradas y capacidad de entrevistar con respeto.
    \item \textbf{Contenido del informe:} hechos, cronología, causas y factores contribuyentes (incluidos los humanos y organizacionales), riesgo, acciones y responsables.
    \item \textbf{Tendencias:} cada trimestre el Gerente SMS revisa los eventos y reportes para buscar repeticiones.
\end{itemize}

\subsection{Consolidación y Registro de Peligros}

Todos los peligros identificados a través de los métodos reactivos, proactivos y predictivos serán registrados y gestionados centralizadamente en el \textbf{“Anexo 7 - Registro de Peligros de Seguridad Operacional"}. Este registro forma parte de nuestro Sistema de Recopilación y Procesamiento de Datos de Seguridad Operacional (SDCPS) y es la base para el siguiente paso del proceso de SRM.

\section{Evaluación y Mitigación de Riesgos de Seguridad Operacional}
Una vez que un peligro ha sido identificado y registrado a través de los métodos descritos en la sección anterior, es imperativo que SKYFILMS lo analice para determinar el riesgo de seguridad operacional asociado. Este proceso de gestión de riesgos (SRM) es sistemático y asegura que se tomen decisiones informadas para mantener los riesgos en un nivel tan bajo como sea razonablemente práctico (ALARP).\\

SKYFILMS aplica dos enfoques de gestión de riesgos, alineados con la naturaleza de la identificación del peligro:

\begin{enumerate}
    \item \textbf{Para Sucesos Reactivos:} Aplicado a los peligros identificados a través de notificaciones de seguridad operacional e investigaciones.
    \item \textbf{Para Operaciones Proactivas y Predictivas:} Aplicado durante la planificación de nuevas operaciones o misiones complejas, utilizando proyecciones estadísticas y auditorias internas y externas.
\end{enumerate}

\subsection{Gestión de Riesgos para Sucesos Reactivos}
Este proceso se activa cuando se recibe una notificación de seguridad (voluntaria u obligatoria) o como resultado de una investigación. La trazabilidad de estos peligros se consigna en el \textbf{Anexo 8 - Matriz de Evaluación de Riesgos}\\

\textbf{a) Índice de Riesgo de Seguridad Operacional:}\\

Para estandarizar la toma de decisiones, todo peligro identificado se evalúa utilizando una matriz de riesgo que cruza la \textbf{probabilidad} de ocurrencia de un suceso con la \textbf{gravedad} de sus consecuencias potenciales.\\

El proceso inicia con la utilización de la matriz de evaluación del riesgo, que asigna un código alfanumérico al riesgo inicial.

\begin{table}[H]
\centering
\renewcommand{\arraystretch}{1.5}
\begin{tabularx}{\textwidth}{|>{\RaggedRight\arraybackslash}p{3.4cm}|*{5}{>{\centering\arraybackslash}X|}}
\hline
\rowcolor[gray]{0.9}\textbf{Probabilidad} & \textbf{Catastrófico (A)} & \textbf{Peligroso (B)} & \textbf{Mayor (C)} & \textbf{Menor (D)} & \textbf{Insignificante (E)} \\ \hline
Frecuente (5) & \cellcolor[HTML]{FF4B4B}\textbf{5A} & \cellcolor[HTML]{FF4B4B}\textbf{5B} & \cellcolor[HTML]{FF4B4B}\textbf{5C} & \cellcolor[HTML]{FFE04D}\textbf{5D} & \cellcolor[HTML]{FFE04D}\textbf{5E} \\ \hline
Ocasional (4) & \cellcolor[HTML]{FF4B4B}\textbf{4A} & \cellcolor[HTML]{FF4B4B}\textbf{4B} & \cellcolor[HTML]{FFE04D}\textbf{4C} & \cellcolor[HTML]{FFE04D}\textbf{4D} & \cellcolor[HTML]{FFE04D}\textbf{4E} \\ \hline
Remoto (3) & \cellcolor[HTML]{FF4B4B}\textbf{3A} & \cellcolor[HTML]{FFE04D}\textbf{3B} & \cellcolor[HTML]{FFE04D}\textbf{3C} & \cellcolor[HTML]{FFE04D}\textbf{3D} & \cellcolor[HTML]{3DBE6B}\textbf{3E} \\ \hline
Improbable (2) & \cellcolor[HTML]{FFE04D}\textbf{2A} & \cellcolor[HTML]{FFE04D}\textbf{2B} & \cellcolor[HTML]{FFE04D}\textbf{2C} & \cellcolor[HTML]{3DBE6B}\textbf{2D} & \cellcolor[HTML]{3DBE6B}\textbf{2E} \\ \hline
Extremadamente improbable (1) & \cellcolor[HTML]{3DBE6B}\textbf{1A} & \cellcolor[HTML]{3DBE6B}\textbf{1B} & \cellcolor[HTML]{3DBE6B}\textbf{1C} & \cellcolor[HTML]{3DBE6B}\textbf{1D} & \cellcolor[HTML]{3DBE6B}\textbf{1E} \\ \hline
\end{tabularx}
\caption{Matriz de evaluación de riesgos: probabilidad (filas) por gravedad (columnas).}
\end{table}


Con el índice de riesgo obtenido, se utiliza la matriz de tolerabilidad para clasificar el riesgo en una de tres regiones:

\begin{itemize}
    \item \textbf{Región NO Tolerable (Roja):} El riesgo es inaceptable y las operaciones relacionadas deben cesar o no iniciarse hasta que el riesgo sea mitigado a un nivel inferior. La gestión es inmediata y requiere la intervención del Gerente SMS y la alta dirección.
    \item \textbf{Región Tolerable (Amarilla):} El riesgo es aceptable siempre que se implementen medidas de mitigación efectivas. Puede requerir una decisión de la dirección para aceptar el riesgo residual.
    \item \textbf{Región Aceptable (Verde):} El riesgo se considera aceptable y puede no requerir acciones de mitigación adicionales, aunque se mantiene bajo observación.
\end{itemize}

\begin{table}[H]
\centering
\renewcommand{\arraystretch}{1.4}
\begin{tabularx}{\textwidth}{|>{\centering\arraybackslash}p{2.6cm}|>{\centering\arraybackslash}p{3.4cm}|>{\RaggedRight\arraybackslash}X|>{\RaggedRight\arraybackslash}p{3.2cm}|}
\hline
\rowcolor[gray]{0.9}\textbf{Región} & \textbf{Índices de riesgo} & \textbf{Criterio} & \textbf{Quién decide y plazo} \\ \hline
\cellcolor[HTML]{FF4B4B}\textbf{No tolerable} & 5A, 5B, 5C, 4A, 4B, 3A & Inaceptable en las condiciones actuales; la operación se detiene. & Gerente General; respuesta inmediata \\ \hline
\cellcolor[HTML]{FFE04D}\textbf{Tolerable} & 5D, 5E, 4C, 4D, 4E, 3B, 3C, 3D, 2A, 2B, 2C & Aceptable si se aplican controles. Requiere decisión de la dirección. & Gerente SMS (ver tabla de niveles de autoridad del capítulo 2) \\ \hline
\cellcolor[HTML]{3DBE6B}\textbf{Aceptable} & 3E, 2D, 2E, 1A, 1B, 1C, 1D, 1E & Aceptable; se mantiene bajo vigilancia. & Gerente SMS; sin plazo \\ \hline
\end{tabularx}
\caption{Matriz de tolerabilidad del riesgo.}
\end{table}


La determinación de la \textbf{probabilidad y gravedad} se basa en las definiciones estandarizadas por SKYFILMS, las cuales han sido adaptadas de la guía de la OACI.

\begin{table}[H]
\centering
\renewcommand{\arraystretch}{1.4}
\begin{tabularx}{\textwidth}{|>{\RaggedRight\arraybackslash}p{4.4cm}|>{\centering\arraybackslash}X|>{\centering\arraybackslash}p{2cm}|}
\hline
\rowcolor[gray]{0.9}\textbf{Probabilidad de ocurrencia} & \textbf{Significado} & \textbf{Valor} \\ \hline
Frecuente & 70 eventos por cada 100 operaciones & \textbf{5} \\ \hline
Ocasional & 50 eventos por cada 100 operaciones & \textbf{4} \\ \hline
Remoto & 25 eventos por cada 100 operaciones & \textbf{3} \\ \hline
Improbable & 10 eventos por cada 100 operaciones & \textbf{2} \\ \hline
Extremadamente improbable & 1 evento por cada 100 operaciones & \textbf{1} \\ \hline
\end{tabularx}
\caption{Probabilidad de ocurrencia.}
\end{table}


\begin{table}[H]
\centering
\renewcommand{\arraystretch}{1.3}
\begin{tabularx}{\textwidth}{|>{\RaggedRight\arraybackslash}p{3cm}|>{\RaggedRight\arraybackslash}X|>{\centering\arraybackslash}p{1.6cm}|}
\hline
\rowcolor[gray]{0.9}\textbf{Gravedad del suceso} & \textbf{Significado} & \textbf{Valor} \\ \hline
Catastrófico & Destrucción del equipo; muertes múltiples. & \textbf{A} \\ \hline
Peligroso & Reducción importante de los márgenes de seguridad, daño físico o una carga de trabajo tal que el personal no puede desempeñar sus tareas de forma precisa y completa; lesiones graves; daños mayores al equipo. & \textbf{B} \\ \hline
Mayor & Reducción significativa de los márgenes de seguridad o de la habilidad del operador para responder a condiciones adversas por aumento de la carga de trabajo o por condiciones que impiden su eficiencia; incidente grave; lesiones a personas. & \textbf{C} \\ \hline
Menor & Interferencia; limitaciones operacionales; uso de procedimientos de emergencia; incidentes menores. & \textbf{D} \\ \hline
Insignificante & Consecuencias leves. & \textbf{E} \\ \hline
\end{tabularx}
\caption{Gravedad de las consecuencias.}
\end{table}


Este flujo de análisis, evaluación y decisión asegura un tratamiento consistente y documentado para cada peligro identificado.

\begin{table}[H]
\centering
\renewcommand{\arraystretch}{1.35}
\begin{tabularx}{\textwidth}{|c|>{\RaggedRight\arraybackslash}p{4.2cm}|>{\RaggedRight\arraybackslash}X|}
\hline
\rowcolor[gray]{0.9}\textbf{Paso} & \textbf{Actividad} & \textbf{Qué ocurre después} \\ \hline
1 & Identificación del peligro & Pasa a evaluar la probabilidad. \\ \hline
2 & Evaluación de la probabilidad (1 a 5) & Pasa a evaluar la gravedad. \\ \hline
3 & Evaluación de la gravedad (A a E) & Con el índice se ubica el riesgo en una región de la matriz de tolerabilidad. \\ \hline
4 & \cellcolor[HTML]{FF4B4B}\textbf{Intolerable} & Acciones correctivas (mitigación TRE); se evalúa de nuevo hasta que el riesgo sea tolerable. \\ \hline
5 & \cellcolor[HTML]{FFE04D}\textbf{Tolerable} & Acciones preventivas o plan de mejora (mitigación TRE); se evalúa de nuevo hasta que el riesgo sea aceptable. \\ \hline
6 & \cellcolor[HTML]{3DBE6B}\textbf{Aceptable} & Se acepta y se vigila. \\ \hline
7 & Fin & Se registra el resultado en el registro de peligros (Anexo 7). \\ \hline
\end{tabularx}
\caption{Proceso de gestión de los riesgos de seguridad operacional.}
\end{table}


\textbf{b) Mitigación y Control de los Riesgos de Seguridad Operacional:}\\

Como primera medida y antes de iniciar el proceso de mitigación de riesgos, en caso de tener mas de un riesgo se debe realizar el siguiente procedimiento de priorización:

\begin{enumerate}
    \item Agrupar los riesgos según su naturaleza (NO tolerable o Tolerable).
    \item Priorizar la mitigación de los riesgos presentes en la sección de \textbf{NO Tolerable}, y dentro de estos priorizar estos de la siguiente manera:
    \begin{enumerate}
        \item Riesgo frente a perdidas humanas. 
        \item Riesgo frente a infraestructura o tecnología.
        \item Riesgo frente al medio ambiente.
    \end{enumerate}
    \item Mitigar los riegos en ubicados en la sección \textbf{Tolerable}.
\end{enumerate}

Para todo riesgo clasificado en la región "NO Tolerable" o "Tolerable", se debe desarrollar e implementar un plan de acción para mitigarlo. Las estrategias de mitigación se centran en:

\begin{itemize}
    \item \textbf{Evitar:} Cancelar o posponer la actividad si el riesgo excede los beneficios de continuar.
    \item \textbf{Reducir:} Implementar acciones para disminuir la probabilidad de ocurrencia o la severidad de las consecuencias.
    \item \textbf{Separar la exposición:} Crear barreras o redundancias para protegerse de los efectos del peligro.
\end{itemize}

Como base principal para las defensas implementadas, SKYFILMS utiliza la metodología \textbf{TRE (Tecnología, Reglamento y Entrenamiento)}, conforme al RAC 100.715:

\begin{itemize}
    \item \textbf{Tecnología:} Uso de sistemas como por ejemplo geocercas, paracaídas balísticos, sistemas de detección y evasión (DAA), o UAS con mayor redundancia.
    \item \textbf{Reglamento:} Creación o modificación de procedimientos internos, políticas y listas de verificación documentadas en los manuales de la empresa.
    \item \textbf{Entrenamiento:} Capacitación específica al personal sobre nuevos procedimientos, simulaciones de emergencia o fortalecimiento de competencias.
\end{itemize}

Una vez aplicadas las medidas de mitigación, se realiza una nueva evaluación para determinar el \textbf{riesgo residual}. Si este se encuentra en la región ”Tolerable” aplicando las medidas propuestas en el análisis o "Aceptable", el plan de acción se considera efectivo y el peligro queda bajo control.\\

\textbf{NOTA:} El responsable de decidir si aceptar un riego ubicado en la sección \textbf{Tolerable} es el Gerente SMS.

\textbf{Tolerabilidad, autoridad y plazos.} La región de cada riesgo decide quién lo acepta y en qué plazo se mitiga; la tabla de \textit{Niveles de Autoridad para Aceptar Riesgos} está en el capítulo 2. En esa tabla, el Gerente General decide en la región roja y el Gerente SMS en la amarilla y la verde.

\textbf{Reglas para que la evaluación sea consistente.}
\begin{itemize}
    \item \textbf{Datos que faltan:} se escribe el supuesto y por qué; si falta información, se usa el peor caso creíble.
    \item \textbf{Dos evaluadores:} todo riesgo amarillo o rojo lo evalúan dos personas (el Gerente SMS y el Jefe de Pilotos o el Técnico UAS) y se justifica cualquier diferencia.
    \item \textbf{Revisión periódica:} se vuelven a evaluar los riesgos cada seis meses y siempre que haya un cambio o un evento.
    \item \textbf{Riesgo agregado:} el SAG revisa cada trimestre si varios riesgos pequeños, juntos, hacen un riesgo mayor (por ejemplo, vuelo nocturno más zona urbana más FPV).
    \item \textbf{Visibilidad:} el Gerente SMS entrega al Gerente General la lista de riesgos amarillos y rojos abiertos cada vez que se reúne el SRB.
    \item \textbf{Controles seguros:} antes de aprobar un control, se pregunta si crea un peligro nuevo (por ejemplo, un cambio de procedimiento que aumente la carga de trabajo).
\end{itemize}

\textbf{Verificación de eficacia de cada control.} Se define para cada control: método (revisar registros, repetir la inspección, observar un vuelo o mirar el SPI), frecuencia (mensual, semestral o anual, según el riesgo), responsable, registro (Anexo 8 y Anexo 11) y criterio de éxito. El resultado vuelve al registro de peligros (Anexo 7). Los controles a cargo de un proveedor se verifican con su informe y con una visita o auditoría anual.
